\section{Methodology}
\label{sec:methodology}

\subsection{Information set and marginal filtering}

Let \(r_{i,t}\) be the daily log return of asset \(i\). All prices must be
finite and positive and dates must be unique and increasing. The final 24
monthly holding periods are locked before fitting. On the preceding prefix,
the implementation compares parsimonious sGARCH, GJR-GARCH, and eGARCH
specifications with Gaussian or skewed/standard Student innovations and
AR(0)/AR(1) means. A candidate is admissible only if fixed-lag Ljung--Box
diagnostics for standardised and squared residuals exceed the pre-specified
one-percent gate. For persistent volatility not absorbed by this grid, a
causal two-component EWMA fallback is selected by Gaussian quasi-likelihood
and is subject to the same diagnostic gate.

Evaluation-period pseudo-observations use fixed training parameters and the
training residual empirical CDF:
\[
 u_{i,t}=\frac{\operatorname{findInterval}(z_{i,t},
 \{z_{i,s}:s\in\mathcal T_{\rm train}\})+1/2}
 {|\mathcal T_{\rm train}|+1}.
\]
Consequently, a future residual cannot change an earlier rank or a fitted
marginal parameter.

\subsection{Fully dynamic neural D-vine}

A D-vine order is chosen on training pseudo-observations by exhaustive search
over the \(7!\) paths, maximising the sum of adjacent absolute Kendall
correlations. A Student-\(t\) backbone retains all six trees:
\[
 c_t(u_1,\ldots,u_d)=
 \prod_{\ell=1}^{d-1}\prod_{k=1}^{d-\ell}
 c_{k,k+\ell\mid k+1:k+\ell-1,t}.
\]
Every one of the \(d(d-1)/2=21\) unconditional or conditional edges has a
neural forecast for its correlation. Conditional pseudo-observations in
higher trees are recursively recomputed using the dynamically predicted
lower-tree copulas, rather than copied from a static vine. Each edge network
uses lagged conditional uniforms, filtered innovations, log volatilities, a
static shrinkage anchor, and a score feature. It is trained by chronological
validation likelihood and initialised at the static correlation. Degrees of
freedom remain fixed per edge for identification; tail dependence still
changes with correlation. The method therefore claims dynamic dependence,
not dynamic family switching.

A held-out comparison showed that a tree-one truncation lost 0.1596 average
log likelihood per observation relative to the full vine (\(z=5.04\)).
Accordingly, the otherwise parsimonious dynamic-first-tree/static-higher-tree
shortcut is rejected for these data.

\subsection{Monthly synthetic paths and staged training}

Training-only realised monthly log returns define empirical marginal
quantiles. Neural-vine uniforms are mapped through these quantiles. A
stationary Gaussian AR copula, estimated and capped on training monthly
returns, propagates marginal serial dependence while neural-vine innovations
retain contemporaneous cross-sectional dependence. The generator must pass
marginal, moving-block correlation, lower-tail co-exceedance, and temporal
fidelity gates before training can begin.

All generated paths are consumed once in pre-training. Historical
fine-tuning uses all causal 30-month context plus 24-month decision
trajectories before the holdout, in balanced sequential passes. The first
30 months seed the recurrent state and are not reward observations. The final
24 holding periods enter neither stage.

\subsection{Portfolio constraints, costs, and objective}

For actor logits \(x_t\), define long and short budgets
\[
 B_L=(G+N)/2,\qquad B_S=(G-N)/2,
\]
where \(N=1\) is net exposure and \(G=1.5\) is the gross cap. The differentiable
two-book action is
\[
 w_t=B_L\operatorname{softmax}(x_t)
     -B_S\operatorname{softmax}(-x_t).
\]
Thus \(\sum_i w_{i,t}=N\), short weights are allowed, and
\(\sum_i|w_{i,t}|\leq G\). The initial allocation is equal weight, so the first
decision is not charged for a fictitious trade from zero holdings.

With asset gross returns \(G_{t+1}\), turnover
\(\mathrm{TO}_t=\|w_t-w_{t-1}\|_1\), annual short-borrow rate \(b_s\), and
cash-borrow rate \(b_c\), net portfolio gross return is
\[
 g^{\rm net}_{t+1}=
 \left[1+w_t^\top(G_{t+1}-1)\right]
 \exp\left\{-\kappa\mathrm{TO}_t-
 \frac{b_s\sum_i(-w_{i,t})_+ + b_c(\sum_iw_{i,t}-1)_+}{12}\right\}.
\]
The same convention is applied to every benchmark.

We optimise terminal-wealth CRRA. Dense rewards are the telescoping increments
\[
 R_t=U_\gamma(W_{t+1}/W_0)-U_\gamma(W_t/W_0)
     -\lambda\widehat{\mathrm{CVaR}}_{0.95,t},
\]
with discount one. Therefore the utility component sums exactly to
\(U_\gamma(W_T/W_0)-U_\gamma(1)\); it is not a sequence of unrelated myopic
utilities.

\subsection{Recurrent TD3 agent}

The observation contains normalised wealth, last monthly returns, monthly
EWMA volatilities, previous CVaR, current holdings, net/gross exposure, and
Kendall/tail features from every vine edge. A genuine 30-month causal burn-in
seeds each LSTM sequence. Daily conditional-mean coefficients are excluded
because applying them to monthly returns would mix horizons.

The actor and two critics use one-layer LSTMs with input and hidden layer
normalisation. Training uses recurrent TD3: clipped double-Q targets, delayed
actor updates, target-policy smoothing, replay, gradient clipping, an initial
random-exploration phase, and deterministic-algorithm controls. The
deterministic actor entropy term is zero by default.

\subsection{Common evaluation and inference}

Strategies provide ex-ante weights only. The common evaluator applies the
same 24 realised asset-return vectors, exposure constraints, turnover,
borrow fees, and financing costs. CAGR and arithmetic excess-return Sharpe
are reported separately. Paired utility differences use Newey--West standard
errors with Holm correction, and a circular moving-block White Reality Check
controls selection over multiple candidate models. Repeated RL seeds measure
algorithm variability; they are not independent market histories.

